\Needspace{6\baselineskip}
\begin{center}\hypertarget{native-229-article166}{}\pdfbookmark[2]{166.}{article-166}166.\end{center}
Für die erste Bestimmung einer bisher völlig unbekannten Bahn aus vier Beobachtungen muss man Orte verwenden, die keine zu große heliozentrische Bewegung umfassen. Sonst fehlten uns die Mittel, um bequem eine erste Näherung zu gewinnen. Das Verfahren, das wir sogleich angeben werden, ist jedoch so vielseitig anwendbar, dass man ohne Bedenken Beobachtungen verwenden kann, die eine heliozentrische Bewegung von $30$ oder $40$ Grad umfassen, sofern nur die Abstände von der Sonne nicht zu sehr voneinander abweichen. Hat man die Wahl, so empfiehlt es sich, die Zeitabstände zwischen der ersten und zweiten, der zweiten und dritten sowie der dritten und vierten Beobachtung annähernd gleich zu wählen. Doch auch hierbei ist keine übertriebene Vorsicht nötig, wie das folgende Beispiel zeigen wird. Dort betragen die Zeitabstände $48$, $55$ und $59$ Tage, und die heliozentrische Bewegung beträgt mehr als $50^0$.

Unsere Lösung verlangt ferner, dass die zweite und dritte Beobachtung vollständig sind. Daher werden die Breiten oder Deklinationen der ersten und letzten Beobachtung \SourceLeaf{230}vernachlässigt. Wir haben oben zwar darauf hingewiesen, dass es im Interesse größerer Genauigkeit meist vorzuziehen ist, die Elemente an die beiden vollständigen äußeren Beobachtungen und an die Längen oder Rektaszensionen der dazwischenliegenden anzupassen. Bei der ersten Bahnbestimmung wird man es jedoch nicht bedauern, auf diesen Vorteil zu verzichten, denn eine möglichst mühelos gewonnene Näherung ist von weit größerer Bedeutung. Der Verlust betrifft vor allem die Länge des Knotens und die Bahnneigung, wirkt sich auf die übrigen Elemente kaum merklich aus und lässt sich später leicht ausgleichen.

Der Kürze halber stellen wir das Verfahren so dar, dass alle Orte auf die Ekliptik bezogen sind. Demnach setzen wir vier Längen und zwei Breiten als gegeben voraus. Da in unseren Formeln auch die Breite der Erde berücksichtigt wird, lassen sie sich ohne Weiteres auf den Fall übertragen, in dem der Äquator als Grundebene gewählt wird, sofern man die Längen und Breiten durch Rektaszensionen und Deklinationen ersetzt.

Für Nutation, Präzession und Parallaxe sowie für die Aberration gilt auch hier alles im vorhergehenden Abschnitt Dargestellte. Sind die ungefähren Abstände von der Erde noch nicht anderweitig bekannt, sodass sich zur Berücksichtigung der Aberration Verfahren I aus Artikel 118 anwenden ließe, so werden die beobachteten Orte zunächst nur von der Aberration der Fixsterne befreit. Die Zeiten werden korrigiert, sobald sich im Laufe der Rechnung die Abstände näherungsweise bestimmen lassen, wie weiter unten deutlicher werden wird.
